% ---
% id: foundations/rem:noise-floor
% kind: rem
% title: "What domination controls, and why it needs a noise floor"
% summary: >-
%   A weight at which a model dominates a lattice family shows that the family
%   adds no support, not how far it lowers an element. An element with no
%   rational direction in its support forces a weight of at least its top
%   eigenvalue over the constant; noise in the model removes that bound.
% prominence: standard
% status: accepted
% uses: [def:domination, lem:domination, lem:rational-obstruction, lem:absorbing-domination, def:exact-reference]
% mentions: []
% assumes: []
% refs:
%   - "nahar2025imperfect, Thm. 3, Eq. (20)"
% provenance: newly-derived
% notes: "No requirement uses this remark. It explains the limits of domination for a declared device weight (spec/future-scope.md section 13) and why a snap is not related to its model by domination."
% lean:
% ---

No requirement of the engine relates a snapped family to its model by
domination: a run evaluates the numbers its key length rests on with the
model in place of the snap where it holds one, and rests on a named premise
elsewhere. The remark says what domination would control if a weight were
declared or established, and why an idealised model could not be related to
a lattice family that way at a small weight. Reductions other than domination
that could reduce an idealised model exactly, and a declared device weight,
are future scope (\repodoc{spec/future-scope.md}~\S13).

\begin{remark}\label{rem:noise-floor}
Let $G$ be a \tref{reference} (Definition~\ref{def:exact-reference}) and $F$
a family of the same shape whose entries are \tref[lattice value]{lattice
values}, both complete at $c$, with $G$ dominating $F$ at a weight $q<1$
(Definition~\ref{def:domination}). The relation bounds every element of $F$
from above by its reference element, $(1-q)F_j\preceq G_j$. What that bound
controls, and what it does not, comes first. Three facts then decide when the
weight can be small, and together they explain when an idealised model whose
elements are rank-deficient dominates no lattice family at a small weight,
what a declaration should do instead, and why that obstruction belongs to
domination and not to every reduction.

\emph{What domination controls.} The relation of
Definition~\ref{def:domination} is the hypothesis under which the reference
is a noise channel followed by $F$, together with classical flags of
weight at most $q_s$ on block $s$ \cite[Thm.~3, Eq.~(20)]{nahar2025imperfect}.
Domination bounds each element of $F$ from above and not from below: $F$
may lower the probability of an outcome, and may remove support, making
impossible an outcome that the model makes possible, at a weight below the
lattice step. For one block of dimension one with $c=1$,
$G_1=2^{-k}/3$ and $G_2=1-2^{-k}/3$, the only lattice value at $k$ in
$[0,2^{-k}/3]$ is $0$, and the family $F_1=0$, $F_2=1$ that rounds $G_1$
down and lets $G_2$ take the rest is dominated at every $q\ge2^{-k}/3$,
since $G_1-(1-q)F_1=G_1\succeq0$ and $G_2-(1-q)F_2=q-2^{-k}/3$, although
outcome~$1$ has become impossible.

\emph{Support.} At a weight below one, an element of $F$ has no support
outside that of its reference element (Lemma~\ref{lem:domination}\,(ii)): for
$x$ in the kernel of $G_j$,
$0\le\langle x,(G_j-(1-q)F_j)x\rangle=-(1-q)\langle x,F_jx\rangle$, so
$F_jx=0$. An outcome that the model makes impossible on some direction stays
exactly impossible there, so an element of $F$ must fit inside the support of
its reference element, and that is what a lattice family can fail to do at a
small weight.

\emph{Rationality.} A lattice family has Gaussian-rational entries in the
basis in which it is declared, so the support of each of its elements is
spanned by Gaussian-rational vectors. A dominated element therefore lies in
the largest rational subspace $\mathcal{R}_j$ of the support of $G_j$
(Lemma~\ref{lem:rational-obstruction}\,(iii)), and completeness turns this
into the lower bound $q\ge\lmax(P_jG_jP_j)/c$, with $P_j$ the projector onto
the orthogonal complement of $\mathcal{R}_j$
(Lemma~\ref{lem:rational-obstruction}\,(iv)). The bound depends on the basis
in which the lattice family is declared and not on the lattice exponent: a
finer lattice still has Gaussian-rational entries, so no refinement of the
lattice lowers it.

\emph{An example at two photons.} Take two orthogonal polarisation modes $H$
and $V$ and the two-photon block in the orthonormal occupation basis
$\ket{2,0},\ket{1,1},\ket{0,2}$. With the diagonal mode created by
$a_D^{\dagger}=(a_H^{\dagger}+a_V^{\dagger})/\sqrt2$, the state with both
photons in that mode is
\begin{equation*}
  \ket{DD}=\tfrac{1}{\sqrt2}\big(a_D^{\dagger}\big)^{2}\ket{0,0}
  =\tfrac12\ket{2,0}+\tfrac{1}{\sqrt2}\ket{1,1}+\tfrac12\ket{0,2},
\end{equation*}
since $(a_H^{\dagger})^{2}\ket{0,0}=\sqrt2\ket{2,0}$ and
$a_H^{\dagger}a_V^{\dagger}\ket{0,0}=\ket{1,1}$. For an ideal, lossless
analyser in the diagonal basis without dark counts, the element of the
outcome in which only the detector of $D$ clicks is $\dyad{DD}$ on this block.
Every non-zero vector on that ray has second coordinate $\sqrt2$ times its
first, so none has Gaussian-rational coordinates, $\mathcal{R}_j=\{0\}$, and
Lemma~\ref{lem:rational-obstruction}\,(v) gives $q\ge\tr\dyad{DD}/c=1/c$:
with $c=1$ no lattice family declared in this basis is dominated by the
model at a weight below one, whatever the lattice exponent. Such irrational
directions are common in photon-number blocks of two or more photons,
because the occupation basis carries factors $1/\sqrt{n}$; the diagonal basis
already produces one.

\emph{A noise floor removes the obstruction.} A physical detector has loss and
dark counts. In a model that includes both in every port, no outcome is
impossible on a photon-number block, and, under conditions that the module
modelling the detector states, every element is positive definite on every
block. Supports are then no obstruction: let every element of $F$ but one
lie between zero and its reference element and let the remaining one, chosen
positive definite in the reference, take what completeness leaves; the
reference then dominates $F$ at a weight of at most $\delta/(\lambda+\delta)$,
where $\delta$ is the norm of the residue the remaining element absorbs and
$\lambda$ the least eigenvalue of the reference element at that outcome, not
of the element of $F$, which the residue raises above it
(Lemma~\ref{lem:absorbing-domination}\,(iv)).
When the other elements stay within a few lattice steps of their reference
elements, $\delta$ is of the order of the lattice step times the number of
outcomes and the dimension, and the weight is of the order of that over
$\lambda$. Loss alone or dark counts alone need not suffice: what matters is
that no element is rank-deficient on a block, and a model in which some
outcome is still impossible on some block keeps the obstruction for that
element.

\emph{Declaring the floor.} A weight relates $F$ to the model $G$ and says
nothing about whether $G$ describes the device; that is the premise of the
declaration's own model. A larger floor makes $\lambda$ larger and the
weight smaller, which is a temptation and not a licence: a floor above the
noise the device has makes the model wrong, an error no weight can show. The
floor to declare is the smallest loss and dark-count rates one is sure the
device has.

\emph{Rank deficiency alone is no obstruction.} A block whose elements are
rank-deficient but have supports spanned by Gaussian-rational vectors may be
dominated by its model at a weight of the order of the lattice step. For one
block of dimension two with $c=1$ and
$G_1=\mathrm{diag}(3/5,0)$, $G_2=\mathrm{diag}(2/5,0)$,
$G_3=\mathrm{diag}(0,1)$, put $a\coloneqq\lfloor2^{k}\cdot3/5\rfloor\,2^{-k}$
and $F_1=\mathrm{diag}(a,0)$, $F_2=\mathrm{diag}(1-a,0)$, $F_3=G_3$. Then
$G_1-(1-q)F_1$ and $G_3-(1-q)F_3$ are positive semidefinite for every
$q\in[0,1)$, and $G_2-(1-q)F_2\succeq0$ exactly when
$q\ge(3/5-a)/(1-a)$, which is below $\tfrac52\,2^{-k}$ because
$0\le3/5-a<2^{-k}$ and $1-a>2/5$, although no element is positive definite.
Where the support of some element contains no non-zero Gaussian-rational
vector, no lattice family is dominated below the bound of
Lemma~\ref{lem:rational-obstruction}\,(v) at any lattice exponent; there it
is the noise floor, not a finer lattice, that makes a small weight possible.

\emph{What the obstruction does not say.} Domination compares the reference
with $F$ element by element on the same block, and the channel it yields acts
on each block as the identity followed by a branching into $F$ and flags. A
unitary applied to both families preserves domination; what fixes the basis
is the lattice family, Gaussian-rational in the basis in which it is
declared. Lemma~\ref{lem:rational-obstruction} bounds the domination of such
a family and nothing else. A reduction that first applies another channel can
map some idealised models exactly onto a family with Gaussian-rational
entries, with no direction to snap at all. An ideal analyser whose analysed
polarisation $\cos\alpha\ket{H}+\sin\alpha\ket{V}$ has an irrational
$\tan\alpha$, so that its one-photon element has no Gaussian-rational
direction in its support, is the image, under one unitary applied to the
input of every outcome, of the same analyser at an angle at which its
elements are lattice values, such as $\alpha=0$. On the block of $m$ photons
in two modes, the isometry from that block into the $m$-fold tensor power of
$\mathbb{C}^{2}$, which maps $\ket{2,0}$, $\ket{1,1}$ and $\ket{0,2}$ to
$\ket{HH}$, $(\ket{HV}+\ket{VH})/\sqrt2$ and $\ket{VV}$, carries $\ket{DD}$ of
the example above to $\tfrac12(\ket{HH}+\ket{HV}+\ket{VH}+\ket{VV})$, which
has rational coordinates, at dimension $2^{m}$ instead of $m+1$. Such
reductions are not claims of this module; an idealised model that dominates
no lattice family at a small weight may still be reducible exactly by them.
\end{remark}
